\chapter{Background}\label{chap:background}

\section{Genomic Functional Annotation}\label{sec:bg-functional-annotation}

Functional annotation assigns biological meaning to genomic regions. In this
thesis, the relevant annotation target is the promoter region, a regulatory DNA
region involved in transcription initiation.

\section{Promoter Regions and Transcription Start Sites}\label{sec:bg-promoters}

A promoter region is commonly described relative to the transcription start site
(TSS). The TSS provides a coordinate reference for extracting fixed-length DNA
windows and comparing promoter sequence patterns across organisms and datasets.

\section{Machine Learning for DNA Sequences}\label{sec:bg-ml-dna}

Machine learning models require DNA sequences to be transformed into numerical
representations. Common strategies include sparse nucleotide encodings, k-mer
frequency vectors, learned embeddings, and physicochemical-property mappings.

\section{Capsule Neural Networks}\label{sec:bg-capsules}

Capsule neural networks represent groups of features as vectors or matrices and
use routing mechanisms to model part-whole relationships. This thesis studies
whether this modeling principle can contribute to promoter sequence
classification.
